\documentclass[11pt,a4paper]{article}

% =============================================================================
% PACKAGES
% =============================================================================
\usepackage[utf8]{inputenc}
\usepackage[T1]{fontenc}
\usepackage{lmodern}
\usepackage{amsmath,amssymb,amsthm}
\usepackage{geometry}
\usepackage{hyperref}
\usepackage{booktabs}
\usepackage{longtable}
\usepackage{array}
\usepackage{enumitem}
\usepackage{xcolor}
\usepackage{fancyhdr}
\usepackage{float}
\usepackage{verbatim}

% =============================================================================
% LAYOUT
% =============================================================================
\geometry{top=2.5cm,bottom=2.5cm,left=2.2cm,right=2.2cm}
\pagestyle{fancy}
\fancyhf{}
\lhead{FIN Theory Documentation}
\rhead{Release 5 Status (2026-03-05)}
\cfoot{\thepage}
\setlength{\headheight}{14pt}

% =============================================================================
% COMMANDS
% =============================================================================
\newcommand{\Kd}{K(d)}
\newcommand{\QW}[1]{\texttt{QW-#1}}
\newcommand{\Verdict}[1]{\texttt{#1}}

% =============================================================================
% THEOREM ENVIRONMENTS
% =============================================================================
\newtheorem{definition}{Definition}
\newtheorem{remark}{Remark}

% =============================================================================
% TITLE
% =============================================================================
\title{\textbf{Fractal Information Nadsoliton (FIN) Theory}\\
\large Canonical Scientific Documentation --- Release 5 Status Update}
\author{Krzysztof Żuchowski}
\date{2026-03-05}

\begin{document}
\maketitle

\begin{center}
\small
Repository: \url{https://github.com/hyconiek/Fractal-Nadsoliton-Theory}

DOI: \url{https://doi.org/10.5281/zenodo.17645737}
\end{center}

\section*{Abstract}
This paper provides a canonical Release 5 status summary of FIN in scientific form.

\textbf{Current core status:}
\begin{itemize}[leftmargin=1.4em]
    \item \textbf{Internal strict closure: achieved} in the audited locked chain.
    \item \textbf{Release 5.1 full-closure readiness: not yet} (foundational gaps remain open/partial).
    \item \textbf{Community-level confirmation still missing:} independent external multiteam audit/replication.
    \item \textbf{GW raw-strain branch:} treated as a late empirical search test for Nadsoliton signatures, not as definitional core of closure.
\end{itemize}

The internal closure statement is anchored in: \QW{2069}, \QW{2070}, \QW{2071}, \QW{2081}, \QW{2097}, \QW{2094}, with post-chain consolidation updates in \QW{2115}, refreshed \QW{2114}, and terminal theorem-chain closure by \QW{2179}, \QW{2180}, and \QW{2181}.
Section~\ref{sec:kernel_to_observable_pipeline} states the deterministic kernel-to-observable mapping used in the strict chain.

\tableofcontents
\newpage

% =============================================================================
\section{Origin}
% =============================================================================
FIN began from a single high-level research objective:
\begin{quote}
Can one frozen structural kernel support the known cross-sector structure (mass hierarchy, flavor structure, and selected bridge observables) under strict no-scan/no-retune rules?
\end{quote}

\subsection{Open Structural Problems in Contemporary Physics}
Despite its empirical success, contemporary fundamental physics relies on a significant number of experimentally fixed parameters whose origin is not explained within the Standard Model itself. These include the hierarchy of fermion masses, the values of coupling constants, and the parameters of the CKM and PMNS mixing matrices.

While the Standard Model accepts these as input parameters, the lack of a generative mechanism suggests that the current description may be an effective field theory emerging from a deeper structure. The \textbf{Fractal Information Nadsoliton (FIN) framework} explores the hypothesis that these arbitrary values may arise from the topological and information-theoretic properties of spacetime itself.

\subsection{Working Hypothesis of the FIN Framework}
The FIN framework explores the hypothesis that certain regularities observed in particle physics and cosmology may admit a unified description in terms of structured information geometry.

In this context, the term \textit{Nadsoliton} is used as a technical label for a self-consistent, non-dispersive configuration of interacting informational degrees of freedom.

The core proposal is that physical laws may emerge from the requirement of self-consistency in a fractal information network, leading to the following working structure:
\begin{enumerate}[leftmargin=1.4em]
    \item \textbf{Information as Substrate}: Treating information processing capacity as the fundamental limit of local interactions.
    \item \textbf{Geometry as Emergence}: Modeling forces and particles as geometric and topological features of this substrate.
    \item \textbf{Fractality}: Assuming scale-invariance in the underlying network structure.
\end{enumerate}

% =============================================================================
\section{Philosophy and Scientific Rule Set}
% =============================================================================
\subsection{Conceptual Inspiration: From Philosophy to Formalism}
\begin{remark}
This section describes the \textit{philosophical and conceptual origins} of the FIN framework. These ideas serve as the heuristic source for the mathematical models developed later, but the validity of the physical results does not depend on the acceptance of these philosophical premises.
\end{remark}

\subsection{Epistemological Separation and Scientific Scope}
\begin{remark}[Scope Separation]
All philosophical, theological, or metaphorical inspirations discussed in the origin layer serve exclusively as heuristic motivation.

\textbf{None of the mathematical definitions, derivations, numerical results, or verification studies depend on these premises.}

All physical claims are evaluated solely on:
\begin{enumerate}[leftmargin=1.4em]
    \item internal mathematical consistency,
    \item reproducibility of numerical procedures,
    \item agreement or consistency with known physical limits.
\end{enumerate}
\end{remark}

\subsection{The Eucharistic Inspiration}
The initial conceptual spark for this framework originated from a reflection on the \textbf{Eucharist} and the theological concept of the Logos ("Word") becoming "Flesh" (Matter). This provided a heuristic model for a phase transition from pure information to geometric structure.

While this inspiration is theological, it translates into a concrete physical question:
\begin{quote}
\textit{Is there a consistent mathematical mechanism by which abstract informational constraints can crystallize into stable, tangible geometric structures?}
\end{quote}

\subsection{The Logos Postulate}
\begin{definition}[Logos Postulate]
A heuristic assumption used for conceptual guidance is that reality may be modeled as a dynamic information process, where existence is defined by participation in interaction (processing).
\end{definition}

This postulate motivates the identification of physical quantities with information-theoretic measures:
\begin{itemize}[leftmargin=1.4em]
    \item \textbf{Space}: Emergent from information correlations,
    \item \textbf{Time}: The arrow of entropy increase (information loss),
    \item \textbf{Matter}: Stable topological defects in the information field.
\end{itemize}

\subsection{Fractal Nature of Reality}
Observing self-similarity across experimentally accessible scales, the framework adopts the working hypothesis of fractal structure.
Mathematically, this is parameterized by:
\begin{equation}
D_f = 4 \ln 2 \approx 2.7726.
\end{equation}

\begin{remark}[Status]
The fractal dimension $D_f$ is a structural parameter introduced at the origin layer. In the strict program, its scientific relevance is accepted only through downstream consistency checks.
\end{remark}

\subsection{The Nadsoliton Concept}
\begin{definition}[Nadsoliton]
A Nadsoliton is a self-consistent, non-dispersive field configuration within the FIN framework. The term is descriptive and does not imply a new fundamental particle beyond modeled degrees of freedom.
\end{definition}

Key conceptual properties:
\begin{itemize}[leftmargin=1.4em]
    \item \textbf{Stability}: balance of nonlinearity and dispersion,
    \item \textbf{Self-Organization}: internal dynamics prevent decay,
    \item \textbf{Multi-Scale Coherence}: repeated structure across scales,
    \item \textbf{Maximum Resonance}: operation near information-processing attractor.
\end{itemize}

\subsection{The 12-Octave Structure}
Initial models used 3 octaves (heuristically linked to 3 generations), but later mathematical consistency required expansion to \textbf{12 octaves}. This is linked in project logic to:
\begin{enumerate}[leftmargin=1.4em]
    \item kissing-number geometry in 3D ($K(3)=12$),
    \item FCC nearest-neighbor coordination (12),
    \item closure behavior in generator studies and octave-space numerics.
\end{enumerate}

\subsection{Strict-Rigor Constraints (Release 5)}
FIN uses a hard separation between conceptual and scientific layers:
\begin{itemize}[leftmargin=1.4em]
    \item \textbf{Conceptual layer}: origin and interpretation,
    \item \textbf{Scientific layer}: derivations, gate logic, data handling, and pass/fail criteria.
\end{itemize}

Release 5 chain enforces:
\begin{enumerate}[leftmargin=1.4em]
    \item frozen kernel vector,
    \item no hidden post-lock scan,
    \item no circular anchor backsolve,
    \item reproducible artifacts with explicit provenance.
\end{enumerate}

% =============================================================================
\section{Formal Kernel Core}
% =============================================================================
\subsection{Canonical Effective Kernel}
The working kernel in the strict chain is:
\begin{equation}
\boxed{\Kd = \frac{\cos(\omega d + \phi)}{1 + \beta d^{\eta}}}
\label{eq:kernel}
\end{equation}
where $d=|o-o'|$ is octave distance.

\subsection{Frozen Release 5 Vector (from \QW{2049})}
\begin{align*}
\omega &= 0.18575,\\
\phi &= 0.16250,\\
\beta &= 1.00000,\\
\eta &= 1.80000.
\end{align*}

\subsection{Operational Meaning}
\begin{itemize}[leftmargin=1.4em]
    \item $\cos(\omega d+\phi)$ encodes oscillatory inter-octave structure,
    \item $(1+\beta d^{\eta})^{-1}$ provides nonlinear damping,
    \item same functional form is reused across audited branches in strict chain.
\end{itemize}

\subsection{Canonical Effective FIN Lagrangian (Strict Internal Layer)}
The strict internal derivation stack uses the following canonical effective Lagrangian template:
\begin{equation}
\mathcal{L}_{\mathrm{FIN}}=
\sum_{o=0}^{11}\left(\frac{1}{2}\partial_\mu\Psi_o^\dagger\partial^\mu\Psi_o - V_\Psi(\Psi_o)\right)
+\frac{1}{2}\partial_\mu\Phi\,\partial^\mu\Phi - V_\Phi(\Phi)-\mathcal{L}_{\mathrm{int}}.
\label{eq:fin_lagrangian_full}
\end{equation}
\begin{equation}
\mathcal{L}_{\mathrm{int}}=
\sum_{o=0}^{11} g_Y(\mathrm{gen}(o))\,|\Phi|^2|\Psi_o|^2
+\frac{1}{2}\sum_{o\neq o'}K_{\mathrm{total}}(o,o')\,\Psi_o^\dagger\Psi_{o'}.
\label{eq:fin_lagrangian_interaction}
\end{equation}
\begin{itemize}[leftmargin=1.4em]
    \item $\Psi_o$: octave-indexed matter modes,
    \item $\Phi$: scalar vacuum/order-parameter mode,
    \item $K_{\mathrm{total}}$: structured mixing operator induced by the frozen kernel branch.
\end{itemize}
\begin{remark}
Equation-level closure of this effective template in strict internal scope is strong; however, full fundamental closure (including global reviewer-facing objections listed later) is not claimed yet.
\end{remark}

% =============================================================================
\section{Kernel-to-Observable Derivation Pipeline}
\label{sec:kernel_to_observable_pipeline}
% =============================================================================
\subsection{Step 1: Locked Kernel Feature Extraction}
The strict chain starts from the frozen kernel and computes a fixed feature vector on locked octave support $D$:
\begin{equation}
\mathbf{f}_{K}
=
\left[
\sum_{d\in D} w_d K(d),\;
\sum_{d\in D} w_d\, d\, K(d),\;
\sum_{d\in D} w_d K(d)^2,\;
\sum_{d\in D} w_d\,\Delta_d K(d)
\right].
\end{equation}
Here $(D,w_d)$ are protocol-locked, not target-specific.

\subsection{Step 2: Sector Operator Construction (No Retune)}
Sector operators are deterministic transforms of $\mathbf{f}_K$:
\begin{equation}
\mathbf{o}_{s} = \mathcal{A}_{s}\!\left(\mathbf{f}_{K};\theta_{\mathrm{lock}}\right),
\qquad
s\in\{\text{mass},\text{flavor},\text{EW/gauge},\text{GR/cosmology}\},
\end{equation}
with one locked parameter state $\theta_{\mathrm{lock}}$ for the audited chain.

\subsection{Step 3: Micro Renormalization Constants}
Micro pointwise derivation provides bridge constants:
\[
Z_{\beta},\;\delta_{\eta}
\]
from the strict \QW{2064} branch (with explicit warning-tracking and later tightening path in \QW{2066}).

\subsection{Step 4: Observable Readout Map}
Physical observables are read out by locked sector maps:
\begin{equation}
\widehat{y}_i
=
\mathcal{R}_i\!\left(\mathbf{o}_s,Z_\beta,\delta_\eta,\mathcal{C}_{\mathrm{SI}}\right),
\end{equation}
where $\mathcal{C}_{\mathrm{SI}}$ are SI definition constants when relevant (\texttt{c\_light}, \texttt{hbar}).
No per-observable post-hoc fit is introduced after protocol lock.

\subsection{Step 5: Strict Gate Validation}
The resulting observables are accepted only through explicit gate chain verdicts (\QW{2069}--\QW{2098}, plus \QW{2115} consolidation branch).

% =============================================================================
\section{What the Theory Generates in Release 5}
% =============================================================================
In the current strict scope, FIN generates in gate-audited form:
\begin{enumerate}[leftmargin=1.5em]
    \item strict SM+GR package map with no direct missing / no strict-unresolved IDs,
    \item full baseline radiative-channel program,
    \item strict closure over tracked missing-14 frontier,
    \item strict CKM CP refinement with deterministic phase scheme,
    \item defect-sweep consistency without critical failures,
    \item full scalar ``theory vs readout'' parameter map from package registry.
\end{enumerate}

% =============================================================================
\section{Release 5 Strict Gate Chain}
% =============================================================================
\subsection{Methodology Used for Closure Assessment}
The closure table below is not a standalone declaration; it is the last layer of a locked methodology:
\begin{enumerate}[leftmargin=1.5em]
    \item \textbf{Source locking}: all closure claims are sourced from explicit artifacts (\QW{2069}, \QW{2070}, \QW{2071}, \QW{2081}, \QW{2097}, \QW{2094}, consolidation \QW{2115}/\QW{2114}, and terminal theorem-chain closure \QW{2179}/\QW{2180}/\QW{2181}).
    \item \textbf{No-retune protocol}: after kernel lock, no sector-specific post-hoc retune is accepted.
    \item \textbf{Status typing}: entries are tagged by strict level (\texttt{strict\_internal\_gate}, \texttt{si\_definition}) and by outcome (\texttt{derived}, \texttt{derived\_with\_warning\_reduced}, \texttt{definition\_constant}).
    \item \textbf{Coverage accounting}: the package coverage report must simultaneously satisfy \texttt{n\_missing = 0} and \texttt{n\_strict\_unresolved = 0} in tracked scope.
    \item \textbf{Independent defect sweep}: \QW{2094} runs a separate consistency sweep (130 checks) to prevent table-only overclaim.
\end{enumerate}

\subsection{Canonical Closure Gates}
\begin{center}
\begin{tabular}{p{2.2cm}p{8.8cm}p{2.6cm}}
\toprule
\textbf{Gate} & \textbf{Verdict} & \textbf{Status} \\
\midrule
\QW{2069} & \Verdict{FULL\_SM\_GR\_DERIVATION\_PACKAGE\_PASS} & PASS \\
\QW{2070} & \Verdict{FULL\_RADIATIVE\_PROGRAM\_PASS} & PASS \\
\QW{2071} & \Verdict{SM\_GR\_FULL\_PRECISION\_CLOSURE\_PASS} & PASS \\
\QW{2081} & \Verdict{MISSING14\_STRICT\_RIGOR\_FRONTIER\_PASS\_ALL\_CLOSED} & PASS \\
\QW{2097} & \Verdict{CKM\_CP\_TARGET\_REFINEMENT\_GATE\_PASS\_STRICT} & PASS \\
\QW{2094} & \Verdict{STRICT\_RIGOR\_DEFECT\_SWEEP\_PASS\_NO\_CRITICAL\_DEFECTS} & PASS \\
\bottomrule
\end{tabular}
\end{center}

\subsection{Terminal Theorem-Chain Closure Gates}
\begin{center}
\begin{tabular}{p{2.2cm}p{8.8cm}p{2.6cm}}
\toprule
\textbf{Gate} & \textbf{Verdict} & \textbf{Status} \\
\midrule
\QW{2179} & \Verdict{L13\_U2B2\_EXACT\_MATCHING\_IDENTITY\_GATE\_PASS\_TERMINAL\_CHAIN\_CLOSED} & PASS \\
\QW{2180} & \Verdict{L14\_V2B2\_EXACT\_ACTION\_IDENTIFICATION\_GATE\_PASS\_TERMINAL\_CHAIN\_CLOSED} & PASS \\
\QW{2181} & \Verdict{DUAL\_TERMINAL\_MATCHING\_CLOSURE\_GATE\_PASS} & PASS \\
\bottomrule
\end{tabular}
\end{center}

\subsection{Core Closure Metrics}
\textbf{Package map (\QW{2069}):}
\begin{align*}
 n_{\text{total}} &= 32, & n_{\text{derived strict internal}} &= 30, \\
 n_{\text{SI definition constants}} &= 2, & n_{\text{missing}} &= 0, \\
 n_{\text{strict unresolved}} &= 0.
\end{align*}

\textbf{Radiative program (\QW{2070}):}
\begin{align*}
 n_{\text{channels total}} &= 7, & n_{\text{implemented}} &= 7, \\
 n_{\text{closure-ready}} &= 7, & n_{\text{missing channels}} &= 0.
\end{align*}

\textbf{Full precision gate (\QW{2071}):} pass flags $=6/6$.

\textbf{Defect sweep (\QW{2094}):}
\begin{align*}
 n_{\text{checks}} &= 130, & n_{\text{failed}} &= 0.
\end{align*}

\subsection{Channel Closure Map (from \QW{2069} Registry)}
\begin{center}
\begin{tabular}{p{5.1cm}p{2.1cm}p{2.1cm}p{3.6cm}}
\toprule
\textbf{Group} & \textbf{Entries} & \textbf{Numeric} & \textbf{Max rel.err [\%]} \\
\midrule
\texttt{fermion\_masses} & 12 & 9 & 34.0134 \\
\texttt{flavor} & 4 & 1 & 2.1660 \\
\texttt{gauge\_and\_electroweak} & 8 & 8 & 2.6524 \\
\texttt{gr\_and\_cosmology} & 5 & 5 & 1.3993 \\
\texttt{project\_specific\_bridge\_quantities} & 3 & 3 & 40.6250 \\
\bottomrule
\end{tabular}
\end{center}

\subsection{Anti-Overclaim and Anti-Artifact Safeguards}
\begin{enumerate}[leftmargin=1.5em]
    \item explicit no-scan/no-retune protocol lock across the chain,
    \item strict unresolved tracking: \texttt{strict\_unresolved\_ids = []} in current package scope,
    \item independent defect sweep (\QW{2094}) with 130 checks and no critical failure,
    \item explicit warning channels preserved instead of hidden deletion (\texttt{derived\_with\_warning\_reduced} entries).
\end{enumerate}

\subsection{Post-Chain Consolidation}
\begin{itemize}[leftmargin=1.4em]
    \item \QW{2115}: \Verdict{GRAVITY\_HIERARCHY\_STRICT\_BRIDGE\_GATE\_PASS} ($7/7$).
    \item refreshed \QW{2114}: \Verdict{REMAINING\_STRICT\_CLOSURE\_ROADMAP\_READY} ($6/6$), with \texttt{strict\_unresolved\_ids = []}.
\end{itemize}

\begin{remark}
Internal strict closure in this document means closure inside the audited locked chain. It is not a substitute for independent external replication.
\end{remark}

% =============================================================================
\section{Scalar Parameters: Theory vs Readout}
% =============================================================================
\subsection{Methodology for ``Theory vs Readout'' Tables}
The scalar tables are generated from the package registry (\QW{2069}) under fixed rules:
\begin{enumerate}[leftmargin=1.5em]
    \item \textbf{Inclusion}: scalar entries only; matrix-valued objects are tracked in dedicated gate reports.
    \item \textbf{Reference coupling}: each parameter is paired with its locked reference value from the registry, if defined.
    \item \textbf{Relative error}: for nonzero reference values,
    \[
    \mathrm{rel\_err\_\%} = 100\cdot \frac{\left|\widehat{y}-y_{\mathrm{ref}}\right|}{\left|y_{\mathrm{ref}}\right|}.
    \]
    \item \textbf{Tolerance semantics}: pass/fail is based on per-entry tolerance metadata in the source report, not on a single global threshold.
    \item \textbf{Warning preservation}: warning-tagged bridge quantities remain explicit in the table (no silent removal).
\end{enumerate}

\subsection{Quick Snapshot}
\begin{table}[H]
\centering
\caption{Release 5 Quick Scalar Snapshot}
\label{tab:quick_theory_vs_readout}
\begin{tabular}{|l|c|c|c|}
\hline
\textbf{Parameter} & \textbf{Theory (pred.)} & \textbf{Readout/ref.} & \textbf{Rel.err [\%]} \\
\hline
\texttt{alpha\_em\_inv\_mz} & 136.108244911 & 137.035999084 & 0.6770 \\
\hline
\texttt{sin2\_theta\_w\_mz} & 0.228998499777 & 0.23122 & 0.9608 \\
\hline
\texttt{alpha\_s\_mz} & 0.115384254263 & 0.1179 & 2.1338 \\
\hline
\texttt{m\_w} & 78.7929676146 & 80.377 & 1.9708 \\
\hline
\texttt{m\_h} & 126.363590676 & 125.25 & 0.8891 \\
\hline
\texttt{m\_top} & 173 & 173 & 0 \\
\hline
\texttt{g\_f} & 1.197316e-05 & 1.166379e-05 & 2.6524 \\
\hline
\texttt{delta\_cp\_ckm} & 1.17400766888 & 1.2 & 2.1660 \\
\hline
\texttt{G\_newton} & 6.674299e-11 & 6.674300e-11 & 1.1184e-05 \\
\hline
\texttt{lambda\_cosmological} & 1.090129e-52 & 1.105600e-52 & 1.3993 \\
\hline
\texttt{h0} & 67.6231065306 & 67.4 & 0.3310 \\
\hline
\end{tabular}
\end{table}

\subsection{Full Scalar Table}
The full scalar registry table is provided in Appendix~\ref{sec:appendix_h_param_table}.

% =============================================================================
\section{Matrix and Non-Scalar Objects}
% =============================================================================
For matrix-valued objects (e.g. \texttt{ckm\_matrix\_abs}, \texttt{pmns\_matrix\_abs}), strict status is tracked in dedicated gate artifacts and cannot be reduced to a single scalar relative error.

% =============================================================================
\section{Empirical Branch (Late Stage): LIGO/Virgo Search Test}
% =============================================================================
This branch is explicitly treated as an empirical search for Nadsoliton-like signatures in detector strain data, performed after internal closure work.

\subsection{Methodological Position}
\begin{itemize}[leftmargin=1.4em]
    \item \QW{2116} method-repair gate: PASS (protocol repair conditions satisfied).
    \item Legacy strict anomaly robustness verdict from \QW{1725}: unchanged (not robust).
\end{itemize}

So this branch is currently methodological-pass but scientific-anomaly-inconclusive.

\subsection{Interpretation}
A positive result in this branch would strengthen empirical relevance, but internal strict closure does not depend on this branch being positive.

% =============================================================================
\section{Current Scientific Claim Boundary}
% =============================================================================
\subsection{Supported Now}
\begin{enumerate}[leftmargin=1.5em]
    \item Internal strict closure of audited Release 5 chain.
    \item Explicit scalar and gate-level package map with no strict-unresolved IDs.
    \item Reproducible protocol path for independent reruns.
\end{enumerate}

\subsection{Not Yet Claimed}
\begin{enumerate}[leftmargin=1.5em]
    \item Final community-confirmed ToE status.
    \item External universal confirmation without independent multiteam audit.
\end{enumerate}

% =============================================================================
\section{Reviewer-Facing Open Gaps (Release 5.1 Readiness Boundary)}
% =============================================================================
\textbf{Important boundary statement:} despite strong strict internal closure, the full ``Release 5.1 complete fundamental closure'' label is not used yet.

\begin{center}
\begin{tabular}{p{2.0cm}p{3.2cm}p{7.2cm}}
\toprule
\textbf{Gap} & \textbf{Status} & \textbf{What remains} \\
\midrule
L1/L2 & PARTIAL & single-fundamental-entity completion and global soliton/topological-protection proof \\
L3/L18/L19 & PARTIAL & full spinor+gauge closure beyond anchored domains and bridge-level partials \\
L4/L16/L23 & PARTIAL/OPEN & full gravity action-level bridge and theorem-level SM+GR reduction \\
L5 & PARTIAL & complete quantization/unitarity/renormalization/causality closure in one global theorem package \\
L6/L7/L20/L21 & PARTIAL & global uniqueness/identifiability and robust derivation-vs-calibration separation \\
L9 & OPEN & one strong new falsifiable prediction beyond the current closed set \\
L10 & OPEN & independent external multiteam replication \\
L11 & OPEN/PARTIAL & Planck-scale first-principles derivation without external anchor dependence \\
L12 & PARTIAL & nonperturbative RG fixed-point proof \\
L13/L14 & CLOSED (strict internal) & terminal theorem chains closed via \QW{2179}/\QW{2180}/\QW{2181} \\
L22 & CLOSED (strict branch rule) & vacuum branch closure resolved in strict branch logic \\
\bottomrule
\end{tabular}
\end{center}

These gaps are tracked operationally in:
\begin{itemize}[leftmargin=1.4em]
    \item \texttt{LISTA\_LUK\_DO\_UZUPELNIENIA\_FIN\_V5.md},
    \item \texttt{RAPORT\_GREP\_LUK\_ALL\_STUDIES\_FIN\_V5\_2026-03-04.md},
    \item \texttt{RAPORT\_STAN\_TEORII\_FIN\_V5\_1\_READINESS\_2026-03-05.md}.
\end{itemize}

% =============================================================================
\section{Where Latest Closures Are Fully Listed}
% =============================================================================
For end-of-document consolidated closure notes, see Appendix G:
\begin{itemize}[leftmargin=1.4em]
    \item Section~\ref{sec:addendum_latest_closure} (Addendum),
    \item Subsection~\ref{sec:g6_release5_update} (compact Release 5 strict-closure tag).
\end{itemize}

% =============================================================================
\appendix

\section{Appendix G: 2026-03-05 Addendum (QW-2063 to QW-2181)}
\label{sec:addendum_latest_closure}

\subsection{G1. Strict-Rigor Snapshot}
\begin{itemize}[leftmargin=1.4em]
    \item \QW{2065}: strict internal first-principles closure PASS ($12/12$).
    \item \QW{2067}: strengthened strict internal closure PASS ($3/3$).
    \item \QW{2068}: explicit SM+GR registry established ($32$ targets).
    \item \QW{2069}: full package PASS ($30/32$ strict-derived + $2$ SI definition constants; missing $0$; strict-unresolved $0$).
    \item \QW{2070}: full radiative program PASS (implemented $7/7$; closure-ready $7/7$).
    \item \QW{2071}: full precision closure PASS ($6/6$).
    \item \QW{2081}: missing-14 strict frontier fully closed.
    \item \QW{2094}: defect sweep PASS (130 checks, 0 failed).
    \item \QW{2179}: L13 terminal theorem-chain closure PASS.
    \item \QW{2180}: L14 terminal theorem-chain closure PASS.
    \item \QW{2181}: dual terminal synchronization closure PASS.
\end{itemize}

\subsection{G2. Scientific Position}
Internal strict closure is achieved in the audited chain, while external independent multiteam replication remains mandatory for community-level final confirmation.

\subsection{G3. Empirical Prediction Package (Operational)}
Locked preregistration branch includes PMNS/cosmology/GW holdout tooling (QW-2076 to QW-2080 family and validators) for external test execution under fixed protocol.

\subsection{G4. GW Holdout Practical Note}
Current holdout branch can be operationally run and audited under locked thresholds, but script-level pass is not equivalent to full independence unless external source and environment independence are guaranteed.

\subsection{G5. Release 5 Compact Closure Tag}
\label{sec:g6_release5_update}
\begin{itemize}[leftmargin=1.4em]
    \item \QW{2069}: \Verdict{FULL\_SM\_GR\_DERIVATION\_PACKAGE\_PASS}
    \item \QW{2070}: \Verdict{FULL\_RADIATIVE\_PROGRAM\_PASS}
    \item \QW{2071}: \Verdict{SM\_GR\_FULL\_PRECISION\_CLOSURE\_PASS}
    \item \QW{2081}: \Verdict{MISSING14\_STRICT\_RIGOR\_FRONTIER\_PASS\_ALL\_CLOSED}
    \item \QW{2097}: \Verdict{CKM\_CP\_TARGET\_REFINEMENT\_GATE\_PASS\_STRICT}
    \item \QW{2094}: \Verdict{STRICT\_RIGOR\_DEFECT\_SWEEP\_PASS\_NO\_CRITICAL\_DEFECTS}
    \item \QW{2179}: \Verdict{L13\_U2B2\_EXACT\_MATCHING\_IDENTITY\_GATE\_PASS\_TERMINAL\_CHAIN\_CLOSED}
    \item \QW{2180}: \Verdict{L14\_V2B2\_EXACT\_ACTION\_IDENTIFICATION\_GATE\_PASS\_TERMINAL\_CHAIN\_CLOSED}
    \item \QW{2181}: \Verdict{DUAL\_TERMINAL\_MATCHING\_CLOSURE\_GATE\_PASS}
\end{itemize}

\section{Appendix H: Release 5 Scalar Parameter Table (Theory vs Readout)}
\label{sec:appendix_h_param_table}

\textbf{Columns:}
\begin{itemize}[leftmargin=1.4em]
    \item Theory = \texttt{predicted\_value}
    \item Readout/ref. = \texttt{reference\_value}
    \item Rel.err = \texttt{rel\_err\_pct}
\end{itemize}

\textbf{Note:} matrix-valued entries (\texttt{ckm\_matrix\_abs}, \texttt{pmns\_matrix\_abs}) are excluded.

\begin{longtable}{|p{3.8cm}|p{3.0cm}|p{3.0cm}|p{2.2cm}|p{3.3cm}|}
\caption{Full Release 5 Scalar Parameter Registry (Theory vs Readout)}\label{tab:full_release5_scalar_registry}\\
\hline
\textbf{Parameter} & \textbf{Theory (pred.)} & \textbf{Readout/ref.} & \textbf{Rel.err [\%]} & \textbf{Status} \\
\hline
\endfirsthead
\hline
\textbf{Parameter} & \textbf{Theory (pred.)} & \textbf{Readout/ref.} & \textbf{Rel.err [\%]} & \textbf{Status} \\
\hline
\endhead
\texttt{alpha\_em\_inv\_mz} & 136.108244911 & 137.035999084 & 0.677014930048 & \texttt{derived} \\
\texttt{sin2\_theta\_w\_mz} & 0.228998499777 & 0.23122 & 0.96077338574 & \texttt{derived} \\
\texttt{m\_top} & 173 & 173 & 0 & \texttt{derived} \\
\texttt{m\_bottom} & 4.07820540272 & 4.18 & 2.4352774469 & \texttt{derived} \\
\texttt{m\_charm} & 1.43521589495 & 1.27 & 13.0091255862 & \texttt{derived} \\
\texttt{m\_tau} & 1.72666557138 & 1.7769 & 2.82708248174 & \texttt{derived} \\
\texttt{m\_muon} & 0.126861131081 & 0.1057 & 20.0199915616 & \texttt{derived} \\
\texttt{m\_electron} & 0.000684808714873 & 0.000511 & 34.0134471377 & \texttt{derived} \\
\texttt{z\_beta} & 114.739579994 & 100 & 14.7395799938 & \texttt{derived\_with\_warning\_reduced} \\
\texttt{delta\_eta} & 1.125 & 0.8 & 40.625 & \texttt{derived\_with\_warning\_reduced} \\
\texttt{gravity\_hierarchy\_beta20} & 8.837545e-41 & 1.000000e-40 & 11.6245525858 & \texttt{derived} \\
\texttt{alpha\_s\_mz} & 0.115384254263 & 0.1179 & 2.13379621433 & \texttt{derived} \\
\texttt{g\_f} & 1.197316e-05 & 1.166379e-05 & 2.65243206822 & \texttt{derived} \\
\texttt{m\_w} & 78.7929676146 & 80.377 & 1.97075330683 & \texttt{derived} \\
\texttt{m\_z} & 89.7345824401 & 91.1876 & 1.59343766026 & \texttt{derived} \\
\texttt{m\_h} & 126.363590676 & 125.25 & 0.889094352451 & \texttt{derived} \\
\texttt{v\_higgs} & 243.017802492 & 246.22 & 1.30054321659 & \texttt{derived} \\
\texttt{m\_up} & 0.001924625 & 0.00216 & 10.8969907407 & \texttt{derived} \\
\texttt{m\_down} & 0.004089828125 & 0.00467 & 12.423380621 & \texttt{derived} \\
\texttt{m\_strange} & 0.0768596124864 & 0.093 & 17.355255391 & \texttt{derived} \\
\texttt{m\_nu1} & 0.017437372283 & N/A & N/A & \texttt{derived} \\
\texttt{m\_nu2} & 0.0194489576105 & N/A & N/A & \texttt{derived} \\
\texttt{m\_nu3} & 0.0531136701061 & N/A & N/A & \texttt{derived} \\
\texttt{delta\_cp\_ckm} & 1.17400766888 & 1.2 & 2.16602759364 & \texttt{derived} \\
\texttt{delta\_cp\_pmns} & 0.0155913898576 & N/A & N/A & \texttt{derived} \\
\texttt{G\_newton} & 6.674299e-11 & 6.674300e-11 & 1.118439e-05 & \texttt{derived} \\
\texttt{c\_light} & 2.997925e+08 & 2.997925e+08 & 0 & \texttt{definition\_constant} \\
\texttt{hbar} & 1.054572e-34 & 1.054572e-34 & 6.127193e-08 & \texttt{definition\_constant} \\
\texttt{lambda\_cosmological} & 1.090129e-52 & 1.105600e-52 & 1.39930382552 & \texttt{derived} \\
\texttt{h0} & 67.6231065306 & 67.4 & 0.331018591428 & \texttt{derived} \\

\end{longtable}

\section{Appendix I: Historical Legacy of the Program (Release 4.4 to Release 5 Status)}
\label{sec:appendix_i_legacy_44}

This appendix summarizes the historical development of the research program from the Release 4.4 era to the present Release 5 status. It replaces raw source embedding with a method-centered narrative. The full archival source snapshot remains available in the repository as:
\texttt{TOE\_FINAL\_DOCUMENTATION\_RELEASE\_4\_4\_LEGACY\_FULL.tex}.

\subsection{A. Origin and Early Research Direction (Pre-Strict Era)}
The 4.4 branch accumulated conceptual and exploratory layers:
\begin{itemize}[leftmargin=1.4em]
    \item information-geometry identity as a unifying ansatz,
    \item kernel-driven coupling hypothesis across sectors,
    \item exploratory mapping between topology, hierarchy, and particle phenomenology.
\end{itemize}
Methodologically, this period mixed formal derivations, heuristic proposals, and calibration-guided tests in one narrative frame.

\subsection{B. Falsification-Driven Correction Cycle}
A key strength of the legacy program was explicit falsification rather than silent deletion:
\begin{itemize}[leftmargin=1.4em]
    \item turbulent-ether interpretation was tested and rejected,
    \item non-robust interpretations were reclassified from ``claim'' to ``open'' or ``legacy'',
    \item methodological risk labels were progressively introduced.
\end{itemize}
This phase established a core principle kept in Release 5 status: preserve negative results as first-class scientific outputs.

\subsection{C. Kernel Formalization and Freeze Discipline}
Across late 4.x, the program migrated from broad exploratory forms to a frozen-kernel protocol:
\begin{itemize}[leftmargin=1.4em]
    \item one shared kernel form constrained cross-sector fitting freedom,
    \item no-scan/no-retune discipline was introduced for strict gates,
    \item report chains began separating ``derived'', ``warning-reduced'', and ``pending-data'' classes.
\end{itemize}
This transition was the methodological bridge from narrative coherence to auditable closure logic.

\subsection{D. Campaign Consolidation: From Partial to Strict Internal Closure}
The 4.9 to 5.0 transition reorganized legacy findings into explicit gate chains:
\begin{itemize}[leftmargin=1.4em]
    \item micro--macro bridge repair and strict pass-path stabilization,
    \item full SM+GR package accounting with explicit unresolved counters,
    \item radiative-channel coverage and defect sweeps under locked protocol.
\end{itemize}
In Release 5 status this yields internal strict closure in audited scope, with remaining requirement of independent external replication.

\subsection{E. Evolution of GW Branch Interpretation}
The GW branch remained an empirical test channel, not the starting axiom of the theory:
\begin{itemize}[leftmargin=1.4em]
    \item legacy branch included ambitious global-correlation interpretations,
    \item strict reruns (including method-repair gates) retained robust local/statistical controls,
    \item non-robust global anomaly language was downgraded where strict calibration did not sustain prior strength.
\end{itemize}
This is consistent with Release 5 policy: observational claims are graded by strict protocol outcome, not by narrative priority.

\subsection{F. Legacy-to-Current Mapping Rule}
For traceable scientific continuity, Release 5 uses the following mapping:
\begin{itemize}[leftmargin=1.4em]
    \item \textbf{Legacy insight} $\rightarrow$ \textbf{strictly typed status} (derived / warning-reduced / pending-data),
    \item \textbf{legacy positive claim} $\rightarrow$ retained only if passing locked gate criteria,
    \item \textbf{legacy conceptual content} $\rightarrow$ preserved as historical motivation, not counted as closure evidence.
\end{itemize}
Therefore, the legacy corpus is preserved historically, while closure statements are anchored exclusively in strict Release 5 gate outputs.

\vspace{1cm}
\begin{center}
\rule{0.5\textwidth}{0.4pt}\\[0.4cm]
Release 5 Status Canonical Scientific Paper --- \today
\end{center}

\end{document}
